\documentclass[11pt]{article}
\usepackage[margin=1in]{geometry}
\usepackage{amsmath,amssymb,amsfonts,mathtools,bm}
\usepackage[T1]{fontenc}
\usepackage[utf8]{inputenc}
\title{Aligned Mathematical Structures}
\author{Source-backed grouped formula sample}
\date{}
\begin{document}
\maketitle
\section{Expressions}
\paragraph{Expression 1.} The following expression is taken from a source-backed formula corpus.

\begin{align*}\begin{array}{ll}D=1:&\chi_y(E)=0\,,\\begin{align*}2mm]D=2:&\chi_y(E)=2(1+10y+y^2)(1-y)^{r-2}\,,\\begin{align*}2mm]D=3:&\chi_y(E)=-\chi(E)y(1+y)(1-y)^{r-3}\,,\\begin{align*}2mm]D=4:&\chi_y(E)=(2+(2r-8-\chi(E))y\\begin{align*}2mm]&\hspace{2cm}+(8r+12-4\chi(E))y^2+(2r-8-\chi(E))y^3+2y^4)(1-y)^{r-4}\,,\\begin{align*}2mm]D=5:&\chi_y(E)=-\chi(E)y(1+y)(1+10y+y^2)(1-y)^{r-5}\,,\end{array} \end{align*}

\paragraph{Expression 2.} The following expression is taken from a source-backed formula corpus.

\begin{align*}\begin{aligned}\mathcal{I}_{\mathrm{D} }(\eta_0, \delta ) - \mathcal{I}_{\mathrm{D} }(\eta, 0 ) & = \mathcal{I}_{\mathrm{D} }(\eta_0, \delta ) - \mathcal{I}_{\mathrm{D} }(\eta, \delta ) + \mathcal{I}_{\mathrm{D} }(\eta, \delta ) - \mathcal{I}_{\mathrm{D} }(\eta, 0 ) \\& \leq \Psi (\eta_0-\eta, 0 ) + \Psi (0, \delta ) = \Psi ( |\eta_0-\eta|, 0 ) + \Psi (0, \delta ).\end{aligned}\end{align*}

\paragraph{Expression 3.} The following expression is taken from a source-backed formula corpus.

\begin{align*}\begin{aligned}\phi_\lambda (v, s^{\prime})&=f (s^{\prime})-\lambda_1 \boldsymbol{d}_{\mathcal{S}_1} (s_1, s_1^{\prime})^{p_1} -\lambda_2 \boldsymbol{d}_{ \mathcal{S}_2} (s_2, s_2^{\prime})^{p_2} \\& \geq B + (B - \lambda_1) \boldsymbol{d}_{\mathcal{S}_1} (s_1, s_1^\prime )^{p_1} + (B - \lambda_2) \boldsymbol{d}_{ \mathcal{S}_2 }(s_2, s_2^\prime )^{p_2} \geq B.\end{aligned}\end{align*}

\paragraph{Expression 4.} The following expression is taken from a source-backed formula corpus.

\begin{align*}\begin{array}{rcl}S_{{\rm IIB}} & = & -T_{M_2}l\int d^2\xi\sqrt{|{\tilde g}_{ij}-\frac{1}{{\tilde k}^2}(\partial_i \varrho+{\tilde B}^{(1)}_i)(\partial_j \varrho+{\tilde B}^{(1)}_j)|}\\& & \\& &-\frac{T_{M_2}l}{2}\int d^2\xi\epsilon^{ij}\left[\tilde{B}^{(2)}{}_{ij}+2\partial_{i}\varrho\ \left({\tilde A}^{(1)}{}_{i}+m\pi\alpha^{\prime} b_{j}\right)\right]\, .\\\end{array}\end{align*}

\paragraph{Expression 5.} The following expression is taken from a source-backed formula corpus.

\begin{align*}\begin{array}{lll}\langle q_{i0} - Q_{i0} \rangle & = & 0 \, , \\&& \\\langle \dot{q}_{i0} \rangle & = & 0 \, , \\&& \\m\omega_{i}\omega_{j} \langle (q_{i0}-Q_{i0})(q_{j0}-Q_{j0}) \rangle & = & k_{B}T \delta_{ij} \, , \\&& \\m\langle \dot{q}_{i0}\dot{q}_{j0} \rangle & = & k_{B}T\delta_{ij} \, , \\&& \\\langle \dot{q}_{i0}(q_{j0} - Q_{j0}) \rangle & = & 0 \, .\end{array}\end{align*}
\end{document}
